\chapter{Function Reference}
\label{ch:funcref}

\ccprod{} has fourteen functions. This chapter documents all fourteen; there
are no others, and there is no facility for adding any.

Function names may be typed in any case. \fml{=sum(a1:a9)} and
\fml{=SUM(A1:A9)} are the same formula.

\section{The functions at a glance}

\begin{cctablethree}{Function}{Arguments}{Returns}{1.35in}{1.0in}{2.45in}
\fn{ABS} & 1 & Absolute value \\
\fn{AND} & 1 or more & \msg{TRUE} if every argument is non-zero \\
\fn{AVERAGE} & 1 or more & Mean of the numeric values \\
\fn{COUNT} & 1 or more & How many numeric values there are \\
\fn{IF} & 2 or 3 & One of two values, by a test \\
\fn{INT} & 1 & Largest whole number not above the value \\
\fn{LEN} & 1 & Number of characters \\
\fn{MAX} & 1 or more & Largest numeric value \\
\fn{MIN} & 1 or more & Smallest numeric value \\
\fn{MOD} & 2 & Remainder \\
\fn{NOT} & 1 & \msg{TRUE} if the argument is zero \\
\fn{OR} & 1 or more & \msg{TRUE} if any argument is non-zero \\
\fn{ROUND} & 1 or 2 & Value rounded to a number of places \\
\fn{SUM} & 1 or more & Total of the numeric values \\
\end{cctablethree}

\section{Ranges}
\index{functions!ranges}

\begin{cccaution}
Only \fn{SUM}, \fn{AVERAGE}, \fn{MIN}, \fn{MAX} and \fn{COUNT} accept a range
as an argument. Every other function given a range returns \errv{\#VALUE!}
(\fn{LEN} returns 0). There is no exception and no warning.
\end{cccaution}

Inside a range, those five functions ignore empty cells and text, and count
\msg{TRUE} and \msg{FALSE} as 1 and 0. An error value anywhere in a range
propagates out as the function's result.

\section{Errors in arguments}

If any argument evaluates to an error, that error becomes the function's result
without the function running --- with one exception. \fn{IF} is exempt, so that
it can choose the branch that is not broken:

\begin{ccexamples}[]{0.44\linewidth}
\exn{=SUM(A1,1/0)}{\#DIV/0!}{the error wins}
\exn{=IF(TRUE,1,1/0)}{1}{the broken branch is not taken}
\end{ccexamples}

\section{Alphabetical reference}

% ---------------------------------------------------------------------------
\ccfunc{ABS}{ABS(number)}
\index{ABS function}

Returns the absolute value of \emph{number} --- the value with any minus sign
removed.

\ccrunin{Arguments}
\begin{ccsteps}
\item \emph{number} --- a number, a boolean, or a reference to a cell holding
one. \textbf{Not a range.}
\end{ccsteps}

\begin{ccexamples}{0.44\linewidth}
\ex{=ABS(7)}{7}
\ex{=ABS(-7)}{7}
\ex{=ABS(-3.25)}{3.25}
\ex{=ABS(0)}{0}
\exn{=ABS(A1:A3)}{\#VALUE!}{ranges are not accepted}
\exn{=ABS("x")}{\#VALUE!}{}
\end{ccexamples}

% ---------------------------------------------------------------------------
\ccfunc{AND}{AND(value1, value2, \ldots)}
\index{AND function}

Returns \msg{TRUE} if every argument is non-zero, \msg{FALSE} otherwise.

\ccrunin{Arguments}
\begin{ccsteps}
\item One or more numbers, booleans or cell references. \textbf{Not ranges.}
\end{ccsteps}

\ccrunin{Notes}
\begin{ccsteps}
\item Every argument is evaluated. There is no short-circuiting: a
\errv{\#DIV/0!} in the fourth argument shows up even if the first is already
\msg{FALSE}.
\item Text in any argument gives \errv{\#VALUE!}.
\end{ccsteps}

\begin{ccexamples}{0.44\linewidth}
\ex{=AND(TRUE,TRUE)}{TRUE}
\ex{=AND(1,1,1)}{TRUE}
\ex{=AND(1,0)}{FALSE}
\exn{=AND(A1>0,A1<100)}{}{the usual use}
\exn{=AND(-1)}{TRUE}{non-zero is true, sign and all}
\end{ccexamples}

% ---------------------------------------------------------------------------
\ccfunc{AVERAGE}{AVERAGE(value1, value2, \ldots)}
\index{AVERAGE function}

Returns the arithmetic mean of the numeric values among the arguments.

\ccrunin{Arguments}
\begin{ccsteps}
\item One or more numbers, booleans, cell references or ranges, freely mixed.
\end{ccsteps}

\ccrunin{Notes}
\begin{ccsteps}
\item \textbf{Empty cells and text are excluded from both the total and the
count.} An average over ten cells of which three hold numbers divides by three.
\item \fn{AVERAGE} of nothing numeric is \errv{\#DIV/0!} --- unlike \fn{MIN} and
\fn{MAX}, which return 0.
\end{ccsteps}

\begin{ccexamples}{0.44\linewidth}
\ex{=AVERAGE(2,4,6)}{4}
\ex{=AVERAGE(B2:B7)}{}
\ex{=AVERAGE(B2:B7,10)}{}
\exn{=AVERAGE(Z1:Z9)}{\#DIV/0!}{all nine cells empty}
\end{ccexamples}

% ---------------------------------------------------------------------------
\ccfunc{COUNT}{COUNT(value1, value2, \ldots)}
\index{COUNT function}

Returns how many of the values are numbers.

\ccrunin{Notes}
\begin{ccsteps}
\item Booleans count. Empty cells and text do not.
\item Ranges are accepted.
\end{ccsteps}

\begin{ccexamples}{0.44\linewidth}
\exn{=COUNT(B1:B10)}{3}{three of the ten hold numbers}
\ex{=COUNT(1,2,3)}{3}
\exn{=COUNT("a","b")}{0}{text is not counted}
\exn{=COUNT(Z1:Z9)}{0}{}
\end{ccexamples}

% ---------------------------------------------------------------------------
\ccfunc{IF}{IF(condition, then) \quad IF(condition, then, else)}
\index{IF function}

Evaluates \emph{condition} and returns \emph{then} if it is true, \emph{else}
if it is not.

\ccrunin{Arguments}
\begin{ccsteps}
\item \emph{condition} --- anything that is not text. \textbf{Any non-zero
value is true}, including a negative one. Text gives \errv{\#VALUE!}.
\item \emph{then} and \emph{else} --- any values at all, of any type. They do
not have to be the same type as one another.
\end{ccsteps}

\ccrunin{Notes}
\begin{ccsteps}
\item If \emph{else} is left out and the condition is false, the result is the
boolean \msg{FALSE}.
\item The branch that is not chosen is still evaluated, but an error in it is
discarded --- \fn{IF} is the one function whose arguments' errors do not
propagate automatically.
\item A range as the condition tests as false.
\end{ccsteps}

\begin{ccexamples}{0.44\linewidth}
\ex{=IF(1>0,"yes","no")}{yes}
\ex{=IF(0,"yes","no")}{no}
\exn{=IF(0,"yes")}{FALSE}{no third argument}
\ex{=IF(B4>B3,2026,2025)}{}
\exn{=IF(-1,"yes","no")}{yes}{non-zero is true}
\exn{=IF("x",1,2)}{\#VALUE!}{text condition}
\end{ccexamples}

% ---------------------------------------------------------------------------
\ccfunc{INT}{INT(number)}
\index{INT function}

Returns the largest whole number that is not greater than \emph{number}.

\ccrunin{Notes}
\begin{ccsteps}
\item This is a \textbf{floor}, not a truncation. It rounds towards minus
infinity, so \fn{INT} of a negative number goes further from zero.
\item Not a range.
\end{ccsteps}

\begin{ccexamples}{0.44\linewidth}
\ex{=INT(2.7)}{2}
\exn{=INT(-2.7)}{-3}{not $-2$}
\ex{=INT(5)}{5}
\ex{=INT(-5)}{-5}
\end{ccexamples}

% ---------------------------------------------------------------------------
\ccfunc{LEN}{LEN(value)}
\index{LEN function}

Returns the number of characters in \emph{value}.

\ccrunin{Notes}
\begin{ccsteps}
\item For a number or a boolean, the value is first written out in General
format and \emph{that} is measured. \fml{=LEN(1/3)} is 10, because one third
displays as \lit{.333333333}.
\item For a range, the result is 0.
\item \fml{=LEN("")} is 0.
\end{ccsteps}

\begin{ccexamples}{0.44\linewidth}
\ex{=LEN("hello")}{5}
\ex{=LEN("")}{0}
\exn{=LEN(1/3)}{10}{the length of .333333333}
\exn{=LEN(TRUE)}{4}{}
\ex{=LEN(A1)}{}
\end{ccexamples}

% ---------------------------------------------------------------------------
\ccfunc{MAX}{MAX(value1, value2, \ldots)}
\index{MAX function}

Returns the largest numeric value among the arguments.

\ccrunin{Notes}
\begin{ccsteps}
\item Empty cells and text are skipped. Booleans count as 1 and 0.
\item Ranges are accepted.
\item \textbf{\fn{MAX} of nothing is 0}, not an error --- which is what Excel
does, and is worth remembering when a column you expected to be full is empty.
\end{ccsteps}

\begin{ccexamples}{0.44\linewidth}
\ex{=MAX(3,9,1)}{9}
\ex{=MAX(B2:B7)}{}
\ex{=MAX(Q1!B2:B4,Q2!B2:B4)}{}
\exn{=MAX(Z1:Z9)}{0}{all empty}
\end{ccexamples}

% ---------------------------------------------------------------------------
\ccfunc{MIN}{MIN(value1, value2, \ldots)}
\index{MIN function}

Returns the smallest numeric value among the arguments. Everything said about
\fn{MAX} applies, including that \fn{MIN} of nothing is 0.

\begin{ccexamples}{0.44\linewidth}
\ex{=MIN(3,9,1)}{1}
\ex{=MIN(B2:B7)}{}
\exn{=MIN(Z1:Z9)}{0}{all empty}
\end{ccexamples}

% ---------------------------------------------------------------------------
\ccfunc{MOD}{MOD(number, divisor)}
\index{MOD function}

Returns the remainder after \emph{number} is divided by \emph{divisor}.

\ccrunin{Notes}
\begin{ccsteps}
\item \textbf{The result takes the sign of the divisor}, not of the number.
This is Excel's rule, and it is not the rule used by C or by most calculators.
\item A divisor of zero gives \errv{\#DIV/0!}.
\item Neither argument may be a range.
\end{ccsteps}

\begin{ccexamples}{0.44\linewidth}
\ex{=MOD(7,3)}{1}
\exn{=MOD(-1,3)}{2}{not $-1$}
\exn{=MOD(1,-3)}{-2}{the divisor's sign}
\ex{=MOD(10,5)}{0}
\ex{=MOD(5,0)}{\#DIV/0!}
\end{ccexamples}

% ---------------------------------------------------------------------------
\ccfunc{NOT}{NOT(value)}
\index{NOT function}

Returns \msg{TRUE} if \emph{value} is zero and \msg{FALSE} if it is anything
else.

\ccrunin{Notes}
\begin{ccsteps}
\item Exactly one argument.
\item Text or a range gives \errv{\#VALUE!}.
\end{ccsteps}

\begin{ccexamples}{0.44\linewidth}
\ex{=NOT(TRUE)}{FALSE}
\ex{=NOT(0)}{TRUE}
\ex{=NOT(5)}{FALSE}
\ex{=NOT(A1>10)}{}
\end{ccexamples}

% ---------------------------------------------------------------------------
\ccfunc{OR}{OR(value1, value2, \ldots)}
\index{OR function}

Returns \msg{TRUE} if any argument is non-zero, \msg{FALSE} if all of them are
zero. Everything said about \fn{AND} applies: no short-circuiting, no ranges,
text gives \errv{\#VALUE!}.

\begin{ccexamples}{0.44\linewidth}
\ex{=OR(FALSE,TRUE)}{TRUE}
\ex{=OR(0,0,0)}{FALSE}
\ex{=OR(A1<0,A1>100)}{}
\end{ccexamples}

% ---------------------------------------------------------------------------
\ccfunc{ROUND}{ROUND(number) \quad ROUND(number, places)}
\index{ROUND function}

Rounds \emph{number} to \emph{places} decimal places.

\ccrunin{Arguments}
\begin{ccsteps}
\item \emph{number} --- a number or a reference. Not a range.
\item \emph{places} --- a whole number from $-9$ to 9. If left out, 0. Anything
outside that range, or anything that is not a whole number, gives
\errv{\#NUM!}.
\end{ccsteps}

\ccrunin{Notes}
\begin{ccsteps}
\item Rounding is \textbf{half away from zero}: 2.5 rounds to 3 and $-2.5$
rounds to $-3$. This is not banker's rounding.
\item A negative \emph{places} rounds to tens, hundreds and so on.
\item \fn{ROUND} changes the \emph{value}. To change only what is displayed you
need a number format, which cannot be set from the keyboard --- see Chapter~\ref{ch:formats}.
\end{ccsteps}

\begin{ccexamples}{0.44\linewidth}
\ex{=ROUND(3.14159,2)}{3.14}
\ex{=ROUND(2.5)}{3}
\exn{=ROUND(-2.5)}{-3}{away from zero}
\ex{=ROUND(1234,-2)}{1200}
\ex{=ROUND(8.675,2)}{8.68}
\exn{=ROUND(1,10)}{\#NUM!}{places out of range}
\end{ccexamples}

% ---------------------------------------------------------------------------
\ccfunc{SUM}{SUM(value1, value2, \ldots)}
\index{SUM function}

Adds up every numeric value among the arguments.

\ccrunin{Arguments}
\begin{ccsteps}
\item One or more numbers, booleans, cell references or ranges, in any mixture.
\end{ccsteps}

\ccrunin{Notes}
\begin{ccsteps}
\item Empty cells and text inside a range are skipped. A bare text argument is
skipped too.
\item Booleans add as 1 and 0.
\item An error anywhere inside a range propagates.
\item \fn{SUM} of nothing numeric is 0.
\item Overflow gives \errv{\#NUM!}.
\end{ccsteps}

\begin{ccexamples}{0.44\linewidth}
\ex{=SUM(1,2,3)}{6}
\ex{=SUM(B2:B7)}{}
\exn{=SUM(B2:B7,C2:C7,100)}{}{mixed freely}
\ex{=SUM(Q1!B2:B40)}{}
\exn{=SUM(Z1:Z9)}{0}{all empty}
\end{ccexamples}

\section{Functions that do not exist}
\index{functions!missing}

For the avoidance of doubt, and because these are the ones people look for
first: there is no \fn{SQRT}, \fn{POWER}, \fn{ROUNDUP}, \fn{ROUNDDOWN},
\fn{TRUNC}, \fn{SIGN}, \fn{SUMIF}, \fn{COUNTIF}, \fn{COUNTA}, \fn{VLOOKUP},
\fn{HLOOKUP}, \fn{INDEX}, \fn{MATCH}, \fn{CONCATENATE}, \fn{LEFT}, \fn{RIGHT},
\fn{MID}, \fn{TRIM}, \fn{UPPER}, \fn{LOWER}, \fn{TEXT}, \fn{VALUE},
\fn{TODAY}, \fn{NOW}, \fn{DATE}, \fn{YEAR}, \fn{MONTH}, \fn{DAY},
\fn{WEEKDAY}, \fn{ISERROR}, \fn{IFERROR}, \fn{ISBLANK}, \fn{MEDIAN},
\fn{STDEV} or \fn{PMT}.

Raising to a power is done with the \lit{\^{}} operator, so \fml{=A1\^{}0.5} is
a square root. \fn{IF} and comparison operators cover a good deal of what
\fn{SUMIF} and \fn{COUNTIF} would be used for, one row at a time.
